\section{Bulk births, deaths and lens exchanges}
\label{sec:peps_bulk_steps}
In the proof of~\cite[Proposition~7.1]{OpenAI2026PolynomialPEPS},
\texttt{06-geometry.tex}, lines 357--445, every bulk birth, death and lens
exchange of the repainting schedule is implemented by the birth lemma
(Theorem~\ref{thm:peps_frame_birth_bounded}) or the exchange lemma
(Theorem~\ref{thm:peps_frame_exchange_bounded}). Both ask for
$I_\Omega(T:E)\le L^{-60}$, and~\cite[Proposition~3.4]{OpenAI2026PolynomialPEPS}
yields this bound from a tapering separation of $T$ from $E$. That
proposition is proved in Theorem~\ref{thm:peps_angular_information_from_gap}.
The bulk changes below are formulated for one slope and threshold; such a pair
can be chosen uniformly from the original spectral gap, above any prescribed
lower bounds. The successive frames and geometric conditions are
required for Proposition~7.1.

Distances are sup distances in the plane. The sites of the open square of side
$L$ are placed at $\Lambda_L=\{1,\dots,L\}^2$, and the capped mark distance
$\delta(x)=\min\{n,d_{\mathcal V}(x)\}$ is that of
Definition~\ref{def:peps_capped_mark_distance}.

\begin{definition}[Information under angular separation]
    \label{def:peps_angular_information_bound}
    \lean{TNLean.PEPS.Approximation.sitePosition,
        TNLean.PEPS.Approximation.AngularInformationBound}
    \leanok
    \uses{def:peps_angular_separation,def:peps_capped_mark_distance}
    Fix $q$, $J$ and $\Delta$. Following~\cite[Proposition~3.4]{OpenAI2026PolynomialPEPS}
    (\texttt{02-information.tex}, lines 635--658) with the separation and the
    diameter bound written pointwise, the statement of information under angular
    separation is the following. For every $a>0$, $A_0\ge1$, integer
    $m_0\ge0$ and $p>0$ there are $D$ and $L_0$ such that for every
    $L\ge L_0$, every Hamiltonian of the model on the open square of side $L$
    with a unit ground vector $\Omega$ and full-system gap $\Delta$, every set
    $\mathcal V$ of at most $m_0$ marked points, every pair of disjoint sets
    $X,E\subseteq\Lambda_L$ and every scale $n$ with
    $D\log L\le n\le A_0L$ and $\operatorname{diam}_\infty X\le A_0n$: if
    $X$ and $E$ satisfy the angular separation condition
    (\ref{eq:peps_angular_separation}) with constant $a$ and floor
    $t=D\log L$, then $I_\Omega(X:E)\le L^{-p}$.
\end{definition}

\begin{definition}[A slope and a threshold at one accuracy]
    \label{def:peps_angular_scale}
    \lean{TNLean.PEPS.Approximation.IsAngularScale}
    \leanok
    \uses{def:peps_angular_information_bound}
    For an integer $p$, a pair $(D,L_0)$ is a slope and a threshold at
    accuracy $L^{-p}$ if the conclusion $I_\Omega(X:E)\le L^{-p}$ of
    Definition~\ref{def:peps_angular_information_bound} holds for every $L\ge L_0$
    and every placement of the sites at a common translation of $\Lambda_L$,
    with the separation and diameter measured in that placement.
\end{definition}

\begin{lemma}[Checking the separation condition]
    \label{lem:peps_angular_separation_checks}
    \lean{TNLean.PEPS.Approximation.cappedMarkDistance_translate,
        TNLean.PEPS.Approximation.AngularSeparation.translate,
        TNLean.PEPS.Approximation.AngularSeparation.union_right,
        TNLean.PEPS.Approximation.AngularSeparation.mono_right,
        TNLean.PEPS.Approximation.AngularSeparation.mono_left,
        TNLean.PEPS.Approximation.AngularSeparation.mono_const,
        TNLean.PEPS.Approximation.angularSeparation_of_frontier,
        TNLean.PEPS.Approximation.min_edgeMarkDist_eq,
        TNLean.PEPS.Approximation.min_cornerMarkDist_eq}
    \leanok
    \uses{def:peps_angular_separation,def:peps_capped_mark_distance,
        def:peps_dyadic_edge_band,lem:peps_dyadic_lens_sides}
    The angular separation condition~(\ref{eq:peps_angular_separation}) is invariant under a common translation
    of $X$, $E$ and $\mathcal V$; it passes to unions and subsets of $E$,
    to subsets of $X$, and, for $t\ge0$, to smaller constants $a'\le a$. If $X\subseteq K$ and $E\subseteq\R^2\setminus K$ (or vice versa)
    and every $x\in X$, $z\in\partial K$ satisfy
    $a\max\{t,\delta(x)\}\le\operatorname{dist}_\infty(x,z)$, then $X$ and $E$ satisfy~(\ref{eq:peps_angular_separation}). The
    capped mark distances $\delta(x)=\min\{n,d_{\mathcal V}(x)\}$ of an edge operation and of the central birth are
    those of the two endpoints of the edge and of the four corners of the
    square.
\end{lemma}

\begin{proof}\leanok
    Translations preserve sup distances. For $x\in X$ and $y\in E$ on opposite sides of $K$ (\texttt{06-geometry.tex}, lines 418--426), the segment $[x,y]$ meets $\partial K$ at some $z$ with
    $\operatorname{dist}_\infty(x,z)\le\operatorname{dist}_\infty(x,y)$. Both definitions of $\delta(x)$ equal $\min\{n,d_{\mathcal V}(x)\}$.
\end{proof}

\begin{theorem}[Slopes and thresholds from information under angular separation]
    \label{thm:peps_angular_scale_exists}
    \lean{TNLean.PEPS.Approximation.AngularInformationBound.exists_isAngularScale}
    \leanok
    \uses{def:peps_angular_information_bound,def:peps_angular_scale,
        lem:peps_angular_separation_checks}
    Assume the statement of
    Definition~\ref{def:peps_angular_information_bound}. For $a>0$,
    $A_0\ge1$, an integer $m_0\ge0$ and an integer $p>0$ there is a slope
    and a threshold $(D,L_0)$ at accuracy $L^{-p}$ in the sense of Definition~\ref{def:peps_angular_scale}.
\end{theorem}

\begin{proof}\leanok
    Taking $(D,L_0)$ from Definition~\ref{def:peps_angular_information_bound} at exponent $p$, translating
    $\mathcal V$ by $-c$ for sites placed at $\Lambda_L+c$ preserves $\operatorname{diam}_\infty X$ and~(\ref{eq:peps_angular_separation}) by Lemma~\ref{lem:peps_angular_separation_checks}.
\end{proof}

\begin{theorem}[Uniform slopes above prescribed lower bounds]
    \label{thm:peps_angular_scale_lower_bounds}
    \lean{TNLean.PEPS.Approximation.exists_isAngularScale_ge}
    \leanok
    \uses{def:peps_angular_scale,def:polynomial_peps_model}
    Fix $q\ge1$, $J\ge0$, $\Delta>0$, $a>0$, $A_0\ge1$, integers
    $m_0\ge0$ and $p>0$, and real lower bounds $D_{\min},L_{\min}$.
    There are $D\ge\max(D_{\min},1)$ and
    $L_0\ge\max(L_{\min},2)$ such that $(D,L_0)$ is a slope and a
    threshold at accuracy $L^{-p}$ (Definition~\ref{def:peps_angular_scale}), chosen before every finite
    system, Hamiltonian, ground vector, translated placement, marked set
    and pair of physical regions.
\end{theorem}
\begin{proof}
    \leanok
    \uses{thm:peps_angular_information_from_gap,thm:peps_angular_scale_exists}
    From the pair $(D_1,L_1)$ provided by Theorem~\ref{thm:peps_angular_information_from_gap} and Theorem~\ref{thm:peps_angular_scale_exists},
    setting $D=\max(D_1,D_{\min},1)$ and
    $L_0=\max(L_1,L_{\min},2)$ gives $\log L\ge0$ and
    $D_1\log L\le D\log L$ for all $L\ge L_0$, which strengthens the floor $t=D\log L$ in~(\ref{eq:peps_angular_separation}) and the lower bound $D\log L\le n$.
\end{proof}

\begin{theorem}[Bulk births, deaths and lens exchanges, conditionally]
    \label{thm:peps_bulk_steps}
    \lean{TNLean.PEPS.Approximation.bulkBirth_step,
        TNLean.PEPS.Approximation.bulkDeath_step,
        TNLean.PEPS.Approximation.lensExchange_step,
        TNLean.PEPS.Approximation.lensExchange_step_of_frontier}
    \leanok
    \uses{def:peps_angular_scale,thm:peps_frame_birth_bounded,
        thm:peps_frame_exchange_bounded,lem:peps_angular_separation_checks}
    Let $(D,L_0)$ be a slope and a threshold at accuracy $L^{-60}$ for
    $(a,A_0,m_0)$, let $L\ge L_0$, and let $\Omega$ be a unit
    gapped ground vector of a Hamiltonian of the model on $\Lambda_L$, with sites placed at a common translation of $\Lambda_L$.
    Let $D\log L\le n\le A_0L$ and $\lvert\mathcal V\rvert\le m_0$.
    \begin{enumerate}
        \item In an encoded frame, let $T$ be a set of physical sites owned by
              $P_\circ$ and disjoint from $H^+(\mathcal F)$, with $\operatorname{diam}_\infty T\le A_0n$, satisfying~(\ref{eq:peps_angular_separation}) with floor
              $t=D\log L$ against the exterior $E$ of the birth partition. Then
              the birth of $T$ at $Q_\circ$ is an allowed word using only
              $P_\circ$ and $Q_\circ$, with at most one pair source and
              reference-vector error at most $L^{-30}$. When these conditions
              hold in the frame after a death, the death is an allowed monomial
              using only $P_\circ$ and $Q_\circ$, with no pair sources, at
              most one pair effect, and error at most $L^{-30}$.
        \item For two encoded frames and the physical sample $Y$ of a lens, with
              every hole outside or inside $Y$, let $(T,E)$ be the exchange partition for a party $C$. If $\operatorname{diam}_\infty T\le A_0n$ and $(T,E)$ satisfies~(\ref{eq:peps_angular_separation}) with $t=D\log L$ (in particular, if $Y$ is the sample of $K\subseteq\R^2$ and $a\max\{t,\delta(x)\}\le\operatorname{dist}_\infty(x,z)$ for all $x\in T$, $z\in\partial K$),
              then the exchange is an allowed word using only $C$, without pair
              resources, with reference-vector error at most $4L^{-30}$.
    \end{enumerate}
\end{theorem}

\begin{proof}\leanok
    As in \texttt{06-geometry.tex} (lines 357--445), Definition~\ref{def:peps_angular_scale} gives $I_\Omega(T:E)\le L^{-60}$,
    from which Theorem~\ref{thm:peps_frame_birth_bounded} and Theorem~\ref{thm:peps_frame_exchange_bounded} yield items~1 and~2; the boundary criterion in item~2 follows from
    Lemma~\ref{lem:peps_angular_separation_checks}.
\end{proof}

\begin{lemma}[The separation premise from the logarithmic floors]
    \label{lem:peps_bulk_separation_from_floors}
    \lean{TNLean.PEPS.Approximation.vertexZone,
        TNLean.PEPS.Approximation.half_angularConstant_le,
        TNLean.PEPS.Approximation.BandOperation.angularSeparation_bandUpdate,
        TNLean.PEPS.Approximation.RepaintingBaseline.angularSeparation_central,
        TNLean.PEPS.Approximation.RepaintingBaseline.floor_exchange_frontier}
    \leanok
    \uses{lem:peps_angular_separation_checks,lem:peps_dyadic_floor,
        lem:peps_dyadic_birth_enlargement,lem:peps_dyadic_hole_enlargement}
    Let $0\le t\le n$, let $a_0$ be the clearance constant of Theorem~\ref{thm:peps_dyadic_angular_separation}, and homogenize the involved guides on the treated squares of
    radius $t$. For an edge birth or death ($\mathcal V$ the two edge endpoints) and for the central birth ($\mathcal V$ the four block corners),
    the closure of the remaining changed region and the union of $\{f_s\ne P_\circ\}$ with the $2\epsilon_0t$-neighbourhood of the true vertices of the surrounding guide $f_s$ satisfy~(\ref{eq:peps_angular_separation}) with constant
    $a_0/2$ and floor $t$ whenever $2\epsilon_0\le a_0/2$. For the lens exchange with $t>0$ and $2\epsilon_0(a_0+1)\le a_0/2$, every point $x$ in the closure of the remaining non-common-$C_e$ positions or within $2\epsilon_0t$ of
    a true vertex of either sheet satisfies
    $(a_0/2)\max\{t,\min(n,d_{\mathcal V}(x))\}\le\operatorname{dist}_\infty(x,z)$
    for all $z\in\partial Y$. All three statements persist under a common translation of
    the sets, $\mathcal V$, and $Y$ (\texttt{06-geometry.tex},
    lines 357--426).
\end{lemma}

\begin{proof}\leanok
    Combining Lemma~\ref{lem:peps_dyadic_floor} (weakened from
    $a_0$ to $a_0/2$) with Lemmas~\ref{lem:peps_dyadic_birth_enlargement} and~\ref{lem:peps_dyadic_hole_enlargement} on each part of the
    exterior and applying Lemma~\ref{lem:peps_angular_separation_checks} gives the bounds and their translation invariance.
\end{proof}

\begin{theorem}[Bulk changes of the repainting from the plane floors, conditionally]
    \label{thm:peps_bulk_steps_from_floors}
    \lean{TNLean.PEPS.Approximation.mem_image_add_right_iff,
        TNLean.PEPS.Approximation.mem_image_vertexZone_of_mem_outerHoles,
        TNLean.PEPS.Approximation.bandBirth_step,
        TNLean.PEPS.Approximation.bandDeath_step,
        TNLean.PEPS.Approximation.centralBirth_step,
        TNLean.PEPS.Approximation.edgeExchange_step}
    \leanok
    \uses{thm:peps_bulk_steps,lem:peps_bulk_separation_from_floors}
    In the setting of Theorem~\ref{thm:peps_bulk_steps}, let $t=D\log L$ with
    $0\le t\le n$, $a\le a_0/2$, and $m_0\ge2$ for edge operations
    ($m_0\ge4$ for the central birth). Place the repainted block at a
    translation $d$, and homogenize the guides on the treated squares of
    radius $t$.
    \begin{enumerate}
        \item For an edge birth or death, or the central birth, with
              $2\epsilon_0\le a_0/2$, let $T$ be physical sites owned by $P_\circ$,
              disjoint from $H^+(\mathcal F)$, with $\operatorname{diam}_\infty T\le A_0n$, placed in the
              translated remaining closed changed region. If sites of other owners lie at positions of labels other than $P_\circ$ in the surrounding guide $f_s$, and
              every hole is centered at a true vertex of $f_s$ with inner
              radius at most $\epsilon_0t$, then the birth or death is the bounded
              change of Theorem~\ref{thm:peps_bulk_steps}(1).
        \item For the lens exchange of an edge, with $t>0$ and
              $2\epsilon_0(a_0+1)\le a_0/2$, let $Y$ be the sites placed in the translated
              lens. If sites not held by $C$ on either sheet lie at
              labels of that sheet's guide other than $C_e$, every hole of either
              sheet is centered at a true vertex of its guide with inner radius
              at most $\epsilon_0t$, and $\operatorname{diam}_\infty T\le A_0n$, then the
              exchange is the bounded change of
              Theorem~\ref{thm:peps_bulk_steps}(2).
    \end{enumerate}
    As in \texttt{06-geometry.tex} (lines 102--108 and 357--445), no separation
    condition remains among the hypotheses.
\end{theorem}

\begin{proof}\leanok
    Every site in an outer hole of inner radius at most $\epsilon_0t$ lies within $2\epsilon_0t$ of its true-vertex center. Thus $T$ lies in the
    translated closed changed region and $E$ lies in the translated union
    of $\{f_s\ne P_\circ\}$ and the $2\epsilon_0t$-neighbourhood of the true vertices (or on opposite sides of the translated
    lens $Y$), so Lemma~\ref{lem:peps_bulk_separation_from_floors} with $a\le a_0/2$ verifies the separation premise of
    Theorem~\ref{thm:peps_bulk_steps}.
\end{proof}

\begin{theorem}[Bulk changes of the repainting for the frames of the sampled guides, conditionally]
    \label{thm:peps_bulk_steps_treated_frames}
    \lean{TNLean.PEPS.Approximation.SquareEdge.endpointBlocks,
        TNLean.PEPS.Approximation.cornerBlocks,
        TNLean.PEPS.Approximation.homogenize_endpointBlocks,
        TNLean.PEPS.Approximation.homogenize_cornerBlocks,
        TNLean.PEPS.Approximation.preimage_edgeTreated,
        TNLean.PEPS.Approximation.preimage_cornerTreated,
        TNLean.PEPS.Approximation.homogenize_preimage_shiftGuide,
        TNLean.PEPS.Approximation.isTrueVertex_shiftGuide_iff,
        TNLean.PEPS.Approximation.IsDisplacedReading.sub_isTrueVertex,
        TNLean.PEPS.Approximation.isTrueVertex_of_mem_treatedTrueVertices,
        TNLean.PEPS.Approximation.guideFrame_holes_isTrueVertex,
        TNLean.PEPS.Approximation.sub_mem_closure_ne_of_reading_ne,
        TNLean.PEPS.Approximation.sub_mem_closure_ne_of_reading_ne_const,
        TNLean.PEPS.Approximation.sub_mem_closure_ne_of_guideFrame_owner_ne,
        TNLean.PEPS.Approximation.lensSample,
        TNLean.PEPS.Approximation.mem_lensSample,
        TNLean.PEPS.Approximation.AngularSeparation.closure_right,
        TNLean.PEPS.Approximation.angularSeparation_birthEnv_guideFrame,
        TNLean.PEPS.Approximation.bandBirth_step_guideFrame,
        TNLean.PEPS.Approximation.bandDeath_step_guideFrame,
        TNLean.PEPS.Approximation.centralBirth_step_guideFrame,
        TNLean.PEPS.Approximation.edgeExchange_step_guideFrame}
    \leanok
    \uses{thm:peps_bulk_steps,lem:peps_bulk_separation_from_floors,
        thm:peps_bulk_steps_from_floors}
    In the setting of Theorem~\ref{thm:peps_bulk_steps_from_floors}, homogenizing
    the guide of the repainted block on the open squares of radius $t$ about the
    edge endpoints or the four block corners equals
    homogenizing it on the treated squares of that operation and translating by $d$. Let each frame be the encoded frame of the sampled guide (Definition~\ref{def:peps_guide_frame}), with site owners read after a common generic
    displacement and holes centered at true vertices of the
    sampled guide with inner radius at most $\epsilon_0t$. Then, after translating back by $-d$:
    (1) every hole center is a true vertex of the homogenized guide;
    (2) every site where the sampled guides before and after a birth or death differ lies in the remaining closed changed region, and the lens sample consists of the sites in $Y+d$; and
    (3) every site not owned by the party of a label $l$ lies in the closure of the positions of labels other than $l$.
    Because~(\ref{eq:peps_angular_separation}) passes to the closure of $E$, the bulk
    changes of Theorem~\ref{thm:peps_bulk_steps_from_floors} hold for these
    frames: for a birth or death on any target $T$ of sites owned
    by $P_\circ$, disjoint from $H^+(\mathcal F)$, with $\operatorname{diam}_\infty T\le A_0n$, where the sampled guides differ; and for the
    exchange on its partition $(T,E)$ with $\operatorname{diam}_\infty T\le A_0n$, without requiring the owner map to be injective (\texttt{06-geometry.tex}, lines 69--80, 102--125, and 328--445).
\end{theorem}

\begin{proof}\leanok
    Translation by $d$ commutes with homogenization on the treated squares and preserves true vertices, giving~(1). By Lemma~\ref{lem:peps_dyadic_sampled_guides}, readings of two guides differ only in the closure of where the guides differ, giving~(2), and a reading takes a value other than $l$ only in the closure of $\{f\ne l\}$, giving~(3). Since the separation inequality in Lemma~\ref{lem:peps_bulk_separation_from_floors} is closed in the exterior variable, it extends to the closure of $E$, and Theorem~\ref{thm:peps_bulk_steps} applies.
\end{proof}

